\section{Lecture 6}

\textbf{Remark}: If \(\frac{d}{dx} F(x, \epsilon)\) exists, then the stochastic Lipschitz condition is equivalent to the derivative condition. Why is this?

Suppose the Lipschitz condition holds. By Fatou's Lemma,
\[\liminf_{x' \to x} E \frac{|F_\epsilon(x') - F_\epsilon(x')|}{|x - x'|} \geq E \liminf_{x' \to x} \frac{|F_\epsilon(x') - F_\epsilon(x')|}{|x - x'|}\]
For the reverse direction, we can use the integral form of the mean-value theorem:
\[\begin{gathered}|F(x, \epsilon) - F(x', \epsilon)| = \int_x^{x'} E \left| \frac{d}{dt} F(t, \epsilon) \right| dt \\
\leq \int_{x'}^x \rho \\
= \rho |x - x'|\end{gathered}\]
\textbf{Remark}: We can generalize to \(p > 1\):
\[\rho_p = \sup_{x \neq x'} \frac{E |F_\epsilon(x) - F_\epsilon(x')^p}{|x - x'|^p} < 1\]
\textbf{Remark}: We can generalize to order 2 autoregressive processes by stacking:
\[Y_i =
\begin{bmatrix}
X_i \\
X_{i-1}
\end{bmatrix}\]

\[Y_i = G(Y_{i-1}, \epsilon_i) \\
= \begin{bmatrix}
F(X_{i-1}, X_i, \epsilon_i) \\
X_{i-1}
\end{bmatrix}\]

For example
\[Y_i =
\begin{bmatrix}
a_1 & a_2 \\
1 & 0
\end{bmatrix}
Y_{i-1}
+
\begin{bmatrix}
\epsilon_i \\
0
\end{bmatrix}\]
The GMC condition reduces to
\[\sup_{y \neq y'}\frac{|A(y - y')|}{|y - y'|} < 1\]
i.e. the operator norm \(|A|_2\) satisfies \(|A|_2 < 1\). However, this operator norm condition is stronger than necessary; the necessary and sufficient condition is that the spectral radius satisfies \(\rho(A) < 1\). Why?
\[\lambda I - A =
\begin{bmatrix}
\lambda - a_1 & -a_2 \\
-1 & \lambda
\end{bmatrix}\]
The characteristic polynomial is
\[\lambda(\lambda - a_1) - a_2 = 0\]
\textbf{Homework}:

\begin{itemize}
\tightlist
\item
  Design an example such that \(\rho(A) < 1\) but \(|A|_2 > 1\)
\item
  Show that \(|A|_2 < 1\) implies that \(\rho(A) < 1\).
\end{itemize}

How can we close the gap between this gap that arises in the vector case, strengthening our result? Look at the two-step ahead iterations
\[\begin{gathered}\sup_{x \neq x'} E \frac{|F_{\epsilon_n} \circ F_{\epsilon_{n-1}}(y) - F_{\epsilon_n} \circ F_{\epsilon_{n-1}}(y')|}{|y - y'|} < 1 \\
= \frac{|A^2 y - A^2 y'|}{|y - y'|}\end{gathered}\]
So we need to impose that \(|A^2|_2 < 1\) in order to assert GMC.

\textbf{Gelfand's Theorem}: For any matrix \(A\),
\[\lim_{n \to \infty} |A^n|_2^{1/n} = \rho(A)\]
and this limiting sequence is monotonically decreasing. This lets us connect the spectral radius theory with the operator norm condition: if \(\rho(A) < 1\), then there exists \(m\) such that \(|A^m|_2 < 1\). Thus, we have
\[\rho_m = \sup_{x \neq x'} \frac{F_{\epsilon_n} \circ \cdots \circ F_{\epsilon_{n-m}}(x) - F_{\epsilon_n} \circ \cdots \circ F_{\epsilon_{n-m}}(x')}{|x - x'|}\]
So \(\frac{\log \rho_m}{m} \to \log \rho(A)\). This is the \emph{Lyapunov exponent} of this dynamical system. How can we estimate \(\rho(A)\) based on data? Or Goldreich is working on this.

\subsection*{Stochastic Gradient Descent}
We have data \(X_1, \dots, X_n\) and a parameter \(\theta\). We want to perform MLE:
\[\max_\theta \sum_{i=1}^n \log \ell(\theta, x_i)\]
Or equivalently
\[\min_\theta \sum_{i=1}^n -\log \ell(\theta, x_i) = \min_\theta f_n(\theta)\]
For modern applications (language modeling training), the sample size \(n\) is extremely large (in the trillions). We could try running gradient descent
\[\theta_j = \theta_{j-1} - \gamma f'(\theta_{j-1})\]
One iteration of gradient descent requires \(O(n)\) computations, competely intractable. \citet{RM51} proposed a separate approach: instead, update with
\[\theta_j = \theta_{j-1} - \gamma_j g(\theta_{j-1}, X_j)\]
reducing the update per step to \(O(1)\).

How to choose the learning rate? Consider the choices \(\gamma_j = \gamma\) or \(\gamma_j = c j^{-\beta}\). This is called \emph{scheduling} in machine learning.

The above update gives us a Markov chain: define
\[F(\theta, x) = \theta - \gamma g(\theta, x)\]
and apply theory from iterated random functions.
\[\begin{gathered}E |F(\theta, x) - F(\theta', x)|^2 \\
= (\theta - \theta')^2 - 2\gamma(\theta - \theta')\E[g(\theta, x) - g(\theta', x)] + \gamma^2 \E[g(\theta, x) - g(\theta', x)]\end{gathered}\]
Define the mean function \(m(\theta) = \E[g(\theta, x)\){]}. Assume that
\[(\theta - \theta')(m(\theta) - m(\theta')) \geq \mu|\theta - \theta')\]
which is \(\mu\)-strong convexity. Next, assume
\[E |g(\theta, X) - g(\theta', X)|^2 \leq L^2 |\theta - \theta'|\]
Then
\[E |F(\theta, X) - F(\theta', X)|^2 \leq (1 - 2\gamma \mu + \gamma^2 L^2)|\theta - \theta'|^2\]
This is a contraction if \(\gamma L < 2\mu\) and \(\gamma < \frac{2\mu}{L^2}\). We can plug in for \(p \neq 2\) and get equivalent conditions for higher moments.
